% FILE: 00_project_identity.tex
% Project: ggen — Governance Generation Engine
% Role in dissertation: manufacturing-machinery-and-code-generation-engine

\section{Project Identity}
\label{sec:ggen-identity}

\subsection{Purpose}
\label{sec:ggen-identity-purpose}

The Governance Generation Engine (ggen) is the manufacturing machinery of the process intelligence
research foundry. Its purpose is to transform ontological knowledge---expressed as RDF triples in
the process intelligence knowledge base---into board-admissible capital-market artifacts and
autonomic runtime code. Where other components in the foundry \emph{produce} process evidence,
ggen \emph{converts} that evidence into artifacts that a board of directors, a financial due
diligence team, or an autonomic runtime can consume.

ggen operates as a configuration-driven pipeline declared in a single TOML manifest
(\texttt{ggen.toml}). Every generation rule in that manifest wires exactly one SPARQL query file
to exactly one Tera template file and exactly one output path. The engine executes each rule,
seals the output with a BLAKE3 cryptographic receipt, and links that receipt to the previous
receipt hash to form an immutable provenance chain. As declared in \texttt{ggen.toml}, the
checksum algorithm is \texttt{blake3} and the proof format is \texttt{cryptographic-chain}.

The project is identified by version \texttt{26.5.21} as recorded in \texttt{audit.json}, with
\texttt{validation\_passed: true} at timestamp \texttt{2026-06-01T20:38:53Z}. This version string
encodes the governance epoch rather than a semantic version, reflecting the foundry's
calendar-based release discipline.

\subsection{Repository Surface}
\label{sec:ggen-identity-surface}

The ggen repository is located at \texttt{/Users/sac/process-intelligence/ggen/}. Its top-level
surface comprises the following directories and files of record:

\begin{itemize}
  \item \texttt{ggen.toml} --- the generation rule manifest; maps SPARQL queries to Tera
    templates and output files for every manufacturing rule.
  \item \texttt{queries/} --- three SPARQL query files: \texttt{extract-board-claims.rq},
    \texttt{extract-diligence-claims.rq}, and \texttt{extract-lifecycle-governance.rq}, plus
    \texttt{select-witness-projections.rq} for the wasm4pm-compat witness lattice.
  \item \texttt{templates/} --- three primary Tera templates: \texttt{ma-deck.tera},
    \texttt{ma-diligence.tera}, and \texttt{blue-river.tera}, plus \texttt{witness-marker.tera}
    and auxiliary projection templates.
  \item \texttt{rules/} --- governance law YAML files: \texttt{feature-law.yaml},
    \texttt{graduation-law.yaml}, \texttt{ts-projection-law.yaml},
    \texttt{wasm-boundary-law.yaml}, and \texttt{component-boundary-law.yaml}.
  \item \texttt{audits/} --- seven audit shell scripts (\texttt{.sh.ggen}) covering 42 or more
    named gates across feature law, TypeScript projection, WASM boundary, and component boundary.
  \item \texttt{ontology-extensions.ttl} and \texttt{wasm4pm-compat.ttl} --- the RDF/Turtle
    ontology extensions defining board claim types and wasm4pm-compat type law.
  \item \texttt{.ggen/receipts/} --- the BLAKE3 cryptographic receipt chain storing four receipt
    JSON files including \texttt{latest.json}.
  \item \texttt{manifests/} --- source specification files including
    \texttt{ts-projection-manifest.yaml} from which projection law was manufactured.
\end{itemize}

Notably, the \texttt{src/} directory exists but contains no Rust source files. ggen is a
specification-and-template engine: it \emph{manufactures} Rust source code as an output artifact
rather than being itself a compiled Rust binary. The \texttt{blue\_river\_dam/src/lib.rs} artifact
is the manufactured Rust output, not ggen's own implementation.

\subsection{Primary Research Role}
\label{sec:ggen-identity-role}

Within the dissertation, ggen occupies the role of \emph{manufacturing machinery and code
generation engine}. It is the component that makes the research foundry's claims capital-market
admissible and runtime-executable. Without ggen, the conformance verdicts produced by wasm4pm
remain laboratory measurements. With ggen, those verdicts are materialized into PowerPoint slides
traceable to BLAKE3 receipts, Excel due-diligence workbooks traceable to SPARQL query results, and
Rust orchestrator source code that encodes the process lifecycle state machine as compilable law.

ggen is designated as a Phase 3 consumer in the foundry's manufacturing sequence, as documented in
\texttt{.ggen/receipts/phase-2-index.json}. It consumes Phase 2 outputs (conformance verdicts,
witness lattice artifacts, lifecycle governance facts) and manufactures Phase 3 capital-market
deliverables. Its position in the pipeline makes it the bridge between the research program and
the board room.

\subsection{Key Artifacts}
\label{sec:ggen-identity-artifacts}

The primary manufactured artifacts attributable to ggen are:

\begin{enumerate}
  \item \textbf{Blue River Orchestrator} (\texttt{../blue\_river\_dam/src/lib.rs}) --- A Rust
    source file manufacturing a MAPE-K autonomic governance engine. The first receipt in the chain
    (\texttt{sync-20260601-175316.json}) records the BLAKE3 output hash of this file as
    \texttt{4e341f5a08c1fd35a5e6aa5dabdd328613263bbc9f7c23bc3bae4893a2562593}.
  \item \textbf{Witness Marker Enumeration}
    (\texttt{../sources/wasm4pm-compat/manufactured/witnesses/witness-markers-20260601.rs}) --- A
    Rust enum defining the eight-position join-semilattice of algebraic proof structures, from
    \texttt{VanDerAalst1989} at the bottom to \texttt{BlueRiverDam} at the top, with embedded
    algebraic test proofs.
  \item \textbf{WIT Boundary Artifact}
    (\texttt{../sources/wasm4pm-compat/manufactured/boundaries/boundary-law-20260601.wit}) --- A
    WebAssembly Interface Types world definition enforcing the compat-world / engine-world
    boundary split.
  \item \textbf{17 governance artifacts} documented in \texttt{GGEN\_MANUFACTURING\_SUMMARY.md}:
    five governance rule files, five Tera templates, and seven audit scripts totaling 42 or more
    named proof gates.
  \item \textbf{BLAKE3 receipt chain} --- three receipts forming an unbroken provenance chain
    from ontology fact to manufactured deliverable, with operation IDs \texttt{de2ad19f},
    \texttt{95a9448c}, and \texttt{09084e1e}.
\end{enumerate}
